%% 1-3 なにも ないときの「0」（10ページ、p.36〜45）
%% 第1巻のきまり：15 pt・行送り28 pt、分かち書き（文節の間に半角の空白）、漢字はすべてルビ
%% 0の たし算・ひき算は 1-4・1-5 で あつかうので、ここでは 出さない。
%% この章で定義した部品（頭文字 bkqc）：\bkqcbasket（たまいれの かご）、\bkqcrocket（ロケット）、\bkqcbadge（A・B の まる）、
%%   \bkqcbutton（エレベーターの ボタン）
\makeatletter
% たまいれの かご  \bkqcbasket{(かごの 下の まん中)}{はいって いる たまの かず（0〜6）}{たまの 色}
\newcommand\bkqcbasket[3]{\begin{scope}[shift={#1}]
  \fill[bkGray!70] (-0.7,-14) rectangle (0.7,0);
  \fill[bkGray!50] (-4,-15) rectangle (4,-13.6);
  \foreach \k [count=\n] in {(-3.3,2.6),(0,2.4),(3.3,2.6),(-1.7,5.4),(1.7,5.4),(0,8.2)} {\ifnum\n>#2\else\fill[#3] \k circle (1.55); \draw[#3!60!black, line width=0.3pt] \k circle (1.55);\fi}
  \draw[bkStem!80, line width=0.5pt] (-6.5,10) -- (-4.2,0) (6.5,10) -- (4.2,0) (-2.2,10) -- (-1.4,0) (2.2,10) -- (1.4,0);
  \draw[bkStem!80, line width=0.5pt] (-5.9,7.3) -- (5.9,7.3) (-5.2,4.4) -- (5.2,4.4) (-4.6,1.8) -- (4.6,1.8);
  \draw[bkStem, line width=1.1pt] (-6.5,10) -- (-4.2,0) -- (4.2,0) -- (6.5,10);
  \draw[bkStem, line width=1.3pt] (0,10) ellipse (6.5 and 1.2);\end{scope}}
% ロケット  \bkqcrocket{(下の まん中)}{大きさ（1で たかさ 約20 mm）}
\newcommand\bkqcrocket[2]{\begin{scope}[shift={#1}, scale=#2]
  \fill[bkOrange] (-1.5,2) .. controls (-2.5,-1.5) and (-1,-2.5) .. (0,-4.5) .. controls (1,-2.5) and (2.5,-1.5) .. (1.5,2) -- cycle;
  \fill[bkYellow] (-0.8,2) .. controls (-1.2,-0.3) and (-0.5,-1) .. (0,-2.4) .. controls (0.5,-1) and (1.2,-0.3) .. (0.8,2) -- cycle;
  \fill[bkRed] (-3,3) -- (-2.6,8) -- (-1.8,9) -- (-1.8,3.5) -- cycle (3,3) -- (2.6,8) -- (1.8,9) -- (1.8,3.5) -- cycle;
  \fill[white] (-2.2,2) rectangle (2.2,14) ;
  \draw[bkGray, line width=0.6pt] (-2.2,2) -- (-2.2,14) (2.2,2) -- (2.2,14);
  \fill[bkRed] (-2.2,14) .. controls (-2.2,17) and (-0.5,19) .. (0,20) .. controls (0.5,19) and (2.2,17) .. (2.2,14) -- cycle;
  \filldraw[fill=bkBlue!50, draw=bkGray, line width=0.6pt] (0,10) circle (1.25);
  \fill[bkRed] (-2.2,2) rectangle (2.2,3.2);\end{scope}}
% A・B の まる  \bkqcbadge{(x,y)}{文字}
\newcommand\bkqcbadge[2]{\node[circle, fill=bkYellow, text=white, font=\sffamily\bfseries\small, minimum size=5.5mm, inner sep=0pt] at #1 {#2};}
% エレベーターの ボタン  \bkqcbutton{(まん中)}{文字}{色}
\newcommand\bkqcbutton[3]{\node[circle, draw=bkGray, line width=0.8pt, fill=#3, minimum size=7mm, inner sep=0pt, font=\sffamily\bfseries\small] at #1 {#2};}
\makeatother

\bkchapter{1}{3}{なにも ないときの「0」}{\bklevel{1}}{なにも ないときの「0」}
\bkmichishirube{1-1,1-2}{1-4,1-6,2-2}

\begin{goal}
\begin{itemize}[label={\color{bkTheme}●}, leftmargin=1.4\zw, itemsep=0pt, topsep=0pt]
\item なにも ない ことを、 0で あらわせる。
\item ものさしの 0に あわせて、 ながさを よめる。
\item 105の ように、 なにも ない くらいに 0を かける。
\end{itemize}
\end{goal}

\begin{hatena}
\begin{bkfigure}
\begin{tikzpicture}[x=1mm, y=1mm]
  % A：リンゴを かぞえる
  \fill[white, rounded corners=3mm] (0,0) rectangle (48,46);
  \bkqcbadge{(6,40)}{A}
  \foreach \i in {1,2,3} {\bkapple{(\i*12,26)}{4.6}{bkAppleRed}
    \node[circle, fill=bkBlue, text=white, font=\sffamily\bfseries\small, minimum size=6.5mm, inner sep=0pt] at (\i*12,13) {\i};}
  \node[font=\sffamily\bfseries\small, text=bkBlue!80!black] at (24,4.5) {1から かぞえる};
  % B：ものさし
  \begin{scope}[shift={(52,0)}]
  \fill[white, rounded corners=3mm] (0,0) rectangle (70,46);
  \bkqcbadge{(6,40)}{B}
  \bkruler{(8,22)}{10.5}{5}
  \draw[bkTheme, line width=1.4pt] (8,16.4) circle (3.6);
  \node[font=\sffamily\bfseries\fontsize{20}{20}\selectfont, text=bkTheme] at (16,36) {？};
  \node[font=\sffamily\bfseries\small, text=bkTheme] at (37,4.5) {0から はじまる};
  \end{scope}
\end{tikzpicture}
\end{bkfigure}
\noindent A：リンゴを かぞえる ときは、 1から はじめるね。\\
B：でも、 ものさしは 0から はじまって いる。\\
どうして 1では ないのかな？
\end{hatena}

\clearpage
%% ---------------- p.37 やってみよう ----------------
\begin{yattemiyou}[ものさしの うえを すすもう]
\begin{enumerate}[label={\bkcircled[bkGreen]{\arabic*}}, leftmargin=2.2\zw, itemsep=2pt, topsep=0pt]
\item けしゴムを、 ものさしの いちばん \R{左}{ひだり}の \R{線}{せん}（\R{目}{め}もり）に おこう。\bkidx{めもり}{\R{目}{め}もり}
\item \R{目}{め}もり 1つ\R{分}{ぶん}ずつ、 \R{右}{みぎ}へ すすめよう。 すすむ たびに、 かずを よもう。
\end{enumerate}
\begin{bkfigure}
\begin{tikzpicture}[x=1mm, y=1mm]
  \fill[white, rounded corners=3mm] (-6.5,-11) rectangle (116.5,27);
  \bkruler{(0,0)}{11}{10}
  \bkhop{(0,0)}{11}{0}{3}{bkGreen}
  \bkeraser{(0,4.4)}{10}{4.2}% 第11回：全巻共通の消しゴム（左はしを 0 の線にあわせて おく）
  \foreach \i/\t in {1/1つ,2/2つ,3/3つ} {\node[font=\sffamily\bfseries\footnotesize, text=bkGreen!70!black] at (\i*11-5.5,14) {\t};}
  \draw[bkBlue!80!black, line width=0.7pt] (1.5,8.9) -- (-3,19);
  \node[anchor=south west, font=\sffamily\bfseries\footnotesize, text=bkBlue!80!black] at (-5.5,19) {けしゴム};
\end{tikzpicture}
\end{bkfigure}
\begin{enumerate}[label={\bkcircled[bkGreen]{\arabic*}}, leftmargin=2.2\zw, itemsep=4pt, topsep=0pt, start=3]
\item \bkkakikomi に かずを かこう。\par\vspace{3pt}
\noindent 1つ すすむと\enspace\bkkakikomi\hspace{2\zw}3つ すすむと\enspace\bkkakikomi\par\vspace{4pt}
\noindent まだ 1つも すすんで いない ときは\enspace\bkkakikomi
\item じぶんの ものさしも みて みよう。 ものさしの はしと 0の \R{目}{め}もりは、 おなじ ところに あるかな？
\end{enumerate}
\end{yattemiyou}

\bkfukidashi{mina}{1つ すすむと 1、 2つ すすむと 2。 じゃあ、 まだ すすんで いない ところは……？}

\clearpage
%% ---------------- p.38〜39 数学の道具 ----------------
\bksection{doguu}{\R{数学}{すう|がく}の \R{道具}{どう|ぐ}}

\bksubsection{1}{なにも ない ときの かず}

\label{bk:1-3:nai}おさらの リンゴを かぞえよう。

\begin{bkfigure}
\begin{tikzpicture}[x=1mm, y=1mm]
  \foreach \px/\n in {0/3,32/2,64/1,96/0} {
    \fill[white, draw=bkLine, line width=0.8pt, rounded corners=2.5mm] (\px,0) rectangle ++(28,22);
    \bkplate{(\px+14,8)}{24}
    \node[font=\sffamily\bfseries\large, text=bkTheme] at (\px+14,-4.5) {\n こ};}
  \foreach \x in {7,14,21} {\bkapple{(\x,11.5)}{3.3}{bkAppleRed}}
  \foreach \x in {42.5,49.5} {\bkapple{(\x,11.5)}{3.3}{bkAppleRed}}
  \bkapple{(78,11.5)}{3.3}{bkAppleRed}
\end{tikzpicture}
\end{bkfigure}

\noindent\begin{minipage}[c]{0.46\linewidth}
いちばん \R{右}{みぎ}の おさらには、 リンゴが ありません。 なにも ない ことも、 かずで 「0」と あらわします。\bkidx{れい}{0（れい）}
\end{minipage}\hfill
\begin{minipage}[c]{0.52\linewidth}\centering
\bkkigocard{0}{れい・ゼロ}{なにも ない}{zero}{1-3}
\end{minipage}\par\vspace{6pt}

\noindent 3こ、 2こ、 1こと、 1こずつ へって いきました。 1より もう 1つ すくないと、 なにも なく なります。 それが 0です。 0は、 1より 1つ すくない かずです。

\noindent どうして 0と かくのかな？ なにも かかないと、 「かぞえるのを わすれたのかな？」と おもわれて しまいます。 「0こ」と かけば、 「かぞえたら 1つも なかった」と はっきり つたわります。

\bksubsection{2}{ものさしの はじまりは 0}

\label{bk:1-3:hajimari}ブロックを ものさしの \R{上}{うえ}に ならべて みよう。\bkidx{ものさし}{ものさし}

\begin{bkfigure}
\begin{tikzpicture}[x=1mm, y=1mm]
  \bkruler{(0,0)}{15}{6}
  \foreach \i/\c in {0/bkRoute2,1/bkRoute4,2/bkRoute6} {\filldraw[fill=\c, draw=\c!60!black, line width=0.6pt] (\i*15,3.3) rectangle ++(15,9);
    \node[font=\sffamily\bfseries\footnotesize, text=white] at (\i*15+7.5,7.8) {\pgfmathparse{int(\i+1)}\pgfmathresult こめ};}
  \draw[bkTheme, line width=1.3pt] (0,-5.6) circle (3);
  \node[anchor=base, font=\sffamily\bfseries\footnotesize, text=bkTheme] at (0,-14.4) {はじまり};
  \draw[bkTheme, line width=1.3pt] (45,-5.6) circle (3);
  \node[anchor=base, font=\sffamily\bfseries\footnotesize, text=bkTheme] at (45,-14.4) {3こ\R{分}{ぶん}};
  \bkbrace{(0,13.5)}{(45,13.5)}{ながさは 3}{1}
\end{tikzpicture}
\end{bkfigure}

\noindent ものさしの 3は、 はじまりから 3こ\R{分}{ぶん} すすんだ ところ。 はじまりは、 まだ なにも すすんで いないので 0です。 これが \bkhatenaref の こたえです。

\noindent もし、 はじまりを 1に したら、 どう なるかな？

\begin{bkfigure}
\begin{tikzpicture}[x=1mm, y=1mm]
  \fill[bkYellowBg, draw=bkYellow!80!black, line width=0.8pt, rounded corners=1mm] (-4,-9) rectangle (94,3);
  \foreach \i in {0,...,6} {\draw[bkInk, line width=0.8pt] (\i*15,3) -- (\i*15,-2.5);
    \pgfmathtruncatemacro\lab{\i+1}\node[font=\sffamily\bfseries\small, text=bkInk] at (\i*15,-5.6) {\lab};}
  \foreach \i/\c in {0/bkRoute2,1/bkRoute4,2/bkRoute6} {\filldraw[fill=\c, draw=\c!60!black, line width=0.6pt] (\i*15,3.3) rectangle ++(15,9);
    \node[font=\sffamily\bfseries\footnotesize, text=white] at (\i*15+7.5,7.8) {\pgfmathparse{int(\i+1)}\pgfmathresult こめ};}
  \draw[bkRed, line width=1.3pt] (45,-5.6) circle (3);
  \node[anchor=west, font=\sffamily\bfseries\small, text=bkRed] at (98,-2) {3こ\R{分}{ぶん}なのに 4？};
  \node[anchor=west] at (98,8) {\bkbatsu};
\end{tikzpicture}
\end{bkfigure}

\noindent はじまりを 1に すると、 ブロック 3こ\R{分}{ぶん}の ところが 4に なって、 ながさと あいません。 はじまりを 0に すると、 すすんだ かずと \R{目}{め}もりの かずが ぴったり おなじに なります。 だから、 ものさしは 0から はじまるのです。

\clearpage
\bksubsection{3}{0に あわせて はかる}

\label{bk:1-3:awaseru}ものさしの はしと 0は、 すこし はなれて います。 はかる ときは、 ものの はしを 0に あわせます。

\begin{bkfigure}
\begin{tikzpicture}[x=1mm, y=1mm]
  \bkruler{(6,0)}{11}{10}
  \bkpencil{(8.2,8.5)}{74.8}{bkRoute6}
  \draw[bkTheme, line width=1pt, dash pattern=on 2pt off 1.5pt] (6,13) -- (6,-2.5);
  \draw[bkTheme, line width=1pt, dash pattern=on 2pt off 1.5pt] (83,13) -- (83,-2.5);
  \draw[bkGray, line width=0.8pt, <-] (2,-9.4) -- (2,-12.5) node[anchor=north, font=\sffamily\bfseries\footnotesize, text=bkGray] {ものさしの はし};
  \node[anchor=south, font=\sffamily\bfseries\footnotesize, text=bkTheme] at (6,13.3) {0に あわせる};
  \node[anchor=south, font=\sffamily\bfseries\footnotesize, text=bkTheme] at (83,13.3) {7で おわる};
\end{tikzpicture}
\end{bkfigure}

\noindent えんぴつの ながさは、 \R{目}{め}もり 7つ\R{分}{ぶん}です。

\noindent どうして はしと 0が はなれて いるのかな？ ものさしの はしは、 おとしたり つかったり すると、 かけたり すりへったり します。 はしを すこし あけて 0の \R{線}{せん}を かいて おけば、 はしが すりへっても ただしく はかれます。

\bksubsection{4}{なにも ない くらいの 0}

\label{bk:1-3:kurai}10の まとまりが 2つで、 ばらが ない とき。 100が 1つ、 10が なくて、 ばらが 5つの とき。

\begin{bkfigure}
\begin{tikzpicture}[x=1mm, y=1mm]
  % 20
  \fill[bkBlue!12, rounded corners=2mm] (0,0) rectangle (22,32);
  \fill[bkOrange!14, rounded corners=2mm] (23,0) rectangle (45,32);
  \node[font=\sffamily\bfseries\fontsize{9}{10}\selectfont, text=bkBlue!80!black] at (11,29) {\R{十}{じゅう}の くらい};
  \node[font=\sffamily\bfseries\fontsize{9}{10}\selectfont, text=bkOrange!80!black] at (34,29) {\R{一}{いち}の くらい};
  \bkrod{(7,9)}{1.6}{bkRoute4}\bkrod{(12,9)}{1.6}{bkRoute4}
  \node[font=\sffamily\bfseries\fontsize{9}{10}\selectfont, text=bkGray] at (34,16) {なし};
  \node[font=\sffamily\bfseries\fontsize{20}{20}\selectfont, text=bkBlue!80!black] at (11,3.6) {2};
  \node[font=\sffamily\bfseries\fontsize{20}{20}\selectfont, text=bkTheme] at (34,3.6) {0};
  % 105
  \begin{scope}[shift={(54,0)}]
  \foreach \px/\c/\t in {0/bkGreen/{\R{百}{ひゃく}の くらい},23/bkBlue/{\R{十}{じゅう}の くらい},46/bkOrange/{\R{一}{いち}の くらい}} {
    \fill[\c!12, rounded corners=2mm] (\px,0) rectangle ++(22,32);
    \node[font=\sffamily\bfseries\fontsize{9}{10}\selectfont, text=\c!80!black] at (\px+11,29) {\t};}
  \bkflat{(3,9)}{1.6}{bkRoute4}
  \node[font=\sffamily\bfseries\fontsize{9}{10}\selectfont, text=bkGray] at (34,16) {なし};
  \foreach \x/\y in {50.5/10,55/10,59.5/10,52.75/14.5,57.25/14.5} {\bkcube{(\x,\y)}{2.2}{bkRoute4}}
  \node[font=\sffamily\bfseries\fontsize{20}{20}\selectfont, text=bkGreen!80!black] at (11,3.6) {1};
  \node[font=\sffamily\bfseries\fontsize{20}{20}\selectfont, text=bkTheme] at (34,3.6) {0};
  \node[font=\sffamily\bfseries\fontsize{20}{20}\selectfont, text=bkOrange!80!black] at (57,3.6) {5};
  \end{scope}
\end{tikzpicture}
\end{bkfigure}

\noindent なにも ない くらいにも 0を かいて、 20、 105と かきます。\bkidx{くらい}{くらい}

\noindent 0を かかないと、 どう なるかな？ 20の 0を かかないと 「2」。 これでは ばらが 2こと おなじに なって しまいます。 105の 0を かかないと 「15」。 10が 1つと ばらが 5こに なって しまいます。 0は、 「この へやは からっぽ」と いう しるし。 0が あるから、 ほかの すうじが ただしい へやに いられるのです。 1-2で みた 10や 100の 0も、 からっぽの へやの 0です。

\begin{taisetsu}
0は、 なにも ない ことを あらわす かず。\\
ものさしの 0は はじまり。 なにも ない くらいにも 0を かく。
\end{taisetsu}

%% （第11回：ここで ページを かえない）
%% ---------------- p.40〜41 例題 ----------------
\begin{reidai}[なにも ない ときの かず]
たまいれの かごの たまは、 なんこ？\par\vspace{3pt}
\noindent\begin{minipage}[c]{0.56\linewidth}
\kangaekata B には 1つも ない。

\kotae A 3こ、\enspace B 0こ
\end{minipage}\hfill
\begin{minipage}[c]{0.42\linewidth}\centering
\begin{tikzpicture}[x=1mm, y=1mm]
  \bkqcbasket{(10,16)}{3}{bkAppleRed}
  \bkqcbasket{(40,16)}{0}{bkAppleRed}
  \bkqcbadge{(0,24)}{A}\bkqcbadge{(30,24)}{B}
\end{tikzpicture}
\end{minipage}
\end{reidai}

\begin{jibun}
おさらの クッキーは なんまい？
\begin{bkfigure}
\begin{tikzpicture}[x=1mm, y=1mm]
  \foreach \px/\l in {0/A,42/B,84/C} {\fill[white, draw=bkLine, line width=0.8pt, rounded corners=2.5mm] (\px,0) rectangle ++(38,16);
    \bkqcbadge{(\px+4.5,11.5)}{\l}
    \bkplate{(\px+21,5.5)}{28}}
  \foreach \x in {16,25} {\bkcookie{(\x,7.5)}{3}}
  \foreach \x in {96,103,110,117} {\bkcookie{(\x,7.5)}{2.9}}
  \foreach \px in {0,42,84} {\node[inner sep=0pt] at (\px+19,-6) {\bkkakikomi\enspace{\small まい}};}
\end{tikzpicture}
\end{bkfigure}
\bkjans{\mbox{A 2まい}\quad \mbox{B 0まい}\quad \mbox{C 4まい}}
\end{jibun}

\begin{reidai}[0までの ならび]
ロケットの はっしゃ！ あいて いる ますに はいる かずは？\par\vspace{3pt}
\noindent\begin{minipage}[c]{0.8\linewidth}
\begin{tikzpicture}[x=1mm, y=1mm, nb/.style={draw=bkBlue, line width=0.8pt, fill=white, rounded corners=1.5mm, minimum width=10mm, minimum height=9mm, font=\sffamily\bfseries}]
  \foreach \x/\t in {0/5,12/4,24/3,36/2,48/1} {\node[nb] at (\x,0) {\t};}
  \node[nb, draw=bkTheme, line width=1.2pt] at (60,0) {0};
\end{tikzpicture}\par\vspace{3pt}
\kangaekata 1こずつ へって いく ならび。 1より 1つ すくないのは 0。\kotaenext 0
\end{minipage}\hfill
\begin{minipage}[c]{0.18\linewidth}\centering
\begin{tikzpicture}[x=1mm, y=1mm]
  \bkqcrocket{(0,0)}{1}
\end{tikzpicture}
\end{minipage}
\end{reidai}

\begin{jibun}
あいて いる ますに はいる かずを かこう。\hspace{1.5\zw}3、\enspace 2、\enspace\bkkakikomi、\enspace 0
\bkjans{1}
\end{jibun}

%% （第11回：ここで ページを かえない）
\begin{reidai}[ものさしで はかる]
テープの ながさは、 \R{目}{め}もり いくつ\R{分}{ぶん}？
\begin{bkfigure}
\begin{tikzpicture}[x=1mm, y=1mm]
  \bkruler{(0,0)}{11}{10}
  \bktape{(0,4)}{66}{bkRoute4}{}
\end{tikzpicture}
\end{bkfigure}
\kangaekata 0から 6まで。\kotaenext 6つ\R{分}{ぶん}
\end{reidai}

\begin{jibun}
えんぴつの ながさは、 \R{目}{め}もり いくつ\R{分}{ぶん}？\par\vspace{3pt}
\noindent\begin{minipage}[c]{0.58\linewidth}\centering
\begin{tikzpicture}[x=1mm, y=1mm]
  \bkruler{(0,0)}{6.5}{10}
  \bkpencil{(2.2,7.5)}{49.8}{bkRoute2}
\end{tikzpicture}
\end{minipage}\hfill
\begin{minipage}[c]{0.4\linewidth}\raggedleft
\R{目}{め}もり\enspace\bkkakikomi\enspace つ\R{分}{ぶん}
\end{minipage}
\bkjans{\R{目}{め}もり 8つ\R{分}{ぶん}}
\end{jibun}

\begin{reidai}[くらいの 0]
100の いたが 1まい、 ばらが 7こ。 いくつ？\par\vspace{3pt}
\noindent\begin{minipage}[c]{0.5\linewidth}
\kangaekata \R{十}{じゅう}の くらいに 0を かく。

\kotae 107
\end{minipage}\hfill
\begin{minipage}[c]{0.48\linewidth}\centering
\begin{tikzpicture}[x=1mm, y=1mm]
  \foreach \px/\c/\t in {0/bkGreen/{\R{百}{ひゃく}の くらい},21/bkBlue/{\R{十}{じゅう}の くらい},42/bkOrange/{\R{一}{いち}の くらい}} {\fill[\c!12, rounded corners=2mm] (\px,0) rectangle ++(20,26);
    \node[font=\sffamily\bfseries\fontsize{8}{9}\selectfont, text=\c!80!black] at (\px+10,23.2) {\t};}
  \bkflat{(2,4)}{1.6}{bkRoute4}
  \foreach \x/\y in {45/5,49.5/5,54/5,58.5/5,47.25/9.5,51.75/9.5,56.25/9.5} {\bkcube{(\x,\y)}{2.2}{bkRoute4}}
  \node[font=\sffamily\bfseries\fontsize{9}{10}\selectfont, text=bkGray] at (31,11) {なし};
\end{tikzpicture}
\end{minipage}
\end{reidai}

\begin{jibun}
100\R{円}{えん}\R{玉}{だま} 3まい、 1\R{円}{えん}\R{玉}{だま} 6まい。 なん\R{円}{えん}？\par\vspace{3pt}
\noindent\begin{minipage}[c]{0.6\linewidth}\centering
\begin{tikzpicture}[x=1.2mm, y=1.2mm]
  \foreach \x in {0,8.5,17} {\bkcoin{(\x,0)}{100}}
  \foreach \x in {28,35,42} {\bkcoin{(\x,3.6)}{1}\bkcoin{(\x,-3.6)}{1}}
\end{tikzpicture}
\end{minipage}\hfill
\begin{minipage}[c]{0.38\linewidth}
\bkkotaeran[6\zw]
\end{minipage}
\bkjans{306\R{円}{えん}}
\end{jibun}

%% ---------------- 第11回：ばらが ない かず・0が 2つ はいる かず ----------------
\begin{reidai}[ばらが ない かず]
10の ぼうが 5\R{本}{ほん}。 ばらは ありません。 いくつ？\par\vspace{3pt}
\noindent\begin{minipage}[c]{0.56\linewidth}
\kangaekata \R{十}{じゅう}の くらいは 5。 \R{一}{いち}の くらいは からっぽなので 0を かく。

\kotae 50
\end{minipage}\hfill
\begin{minipage}[c]{0.42\linewidth}\centering
\begin{tikzpicture}[x=1mm, y=1mm]
  \fill[bkBlue!12, rounded corners=2mm] (0,0) rectangle (25,27);
  \fill[bkOrange!14, rounded corners=2mm] (26,0) rectangle (51,27);
  \foreach \k in {0,...,4} {\bkrod{(2.5+\k*4.2,7)}{1.7}{bkRoute4}}
  \node[font=\sffamily\bfseries\fontsize{9}{10}\selectfont, text=bkGray] at (38.5,15) {なし};
  \node[font=\sffamily\bfseries\large, text=bkBlue!80!black] at (12.5,3) {5};
  \node[font=\sffamily\bfseries\large, text=bkTheme] at (38.5,3) {0};
\end{tikzpicture}
\end{minipage}
\end{reidai}

\begin{jibun}
10\R{円}{えん}\R{玉}{だま}が 7まい。 1\R{円}{えん}\R{玉}{だま}は ありません。 なん\R{円}{えん}？\par\vspace{3pt}
\noindent\begin{minipage}[c]{0.6\linewidth}\centering
\begin{tikzpicture}[x=1mm, y=1mm]
  \foreach \x in {0,8.5,17,25.5,34,42.5,51} {\bkcoin{(\x,0)}{10}}
\end{tikzpicture}
\end{minipage}\hfill
\begin{minipage}[c]{0.38\linewidth}
\bkkotaeran[6\zw]
\end{minipage}
\bkjans{70\R{円}{えん}}
\end{jibun}

\begin{jibun}
100\R{円}{えん}\R{玉}{だま} 1まいと、 10\R{円}{えん}\R{玉}{だま} 6まい。 1\R{円}{えん}\R{玉}{だま}は ありません。 なん\R{円}{えん}？\par\vspace{3pt}
\noindent\begin{minipage}[c]{0.6\linewidth}\centering
\begin{tikzpicture}[x=1mm, y=1mm]
  \bkcoin{(0,0)}{100}
  \foreach \x in {11,19.5,28,36.5,45,53.5} {\bkcoin{(\x,0)}{10}}
\end{tikzpicture}
\end{minipage}\hfill
\begin{minipage}[c]{0.38\linewidth}
\bkkotaeran[6\zw]
\end{minipage}
\bkjans{160\R{円}{えん}}
\end{jibun}

\begin{reidai}[0が 2つ はいる かず]
1000\R{円}{えん}さつ 2まいと、 1\R{円}{えん}\R{玉}{だま} 7まい。 なん\R{円}{えん}？ \R{千}{せん}の くらい（1-2）も つかって かんがえよう。\par\vspace{3pt}
\kangaekata \R{千}{せん}の くらいは 2。 \R{百}{ひゃく}と \R{十}{じゅう}の くらいは からっぽなので 0。 \R{一}{いち}の くらいは 7。
\begin{bkfigure}
\begin{tikzpicture}[x=1mm, y=1mm]
  \foreach \px/\c/\t/\d/\tc in {0/bkPurple/{\R{千}{せん}の くらい}/2/bkPurple,30/bkGreen/{\R{百}{ひゃく}の くらい}/0/bkTheme,60/bkBlue/{\R{十}{じゅう}の くらい}/0/bkTheme,90/bkOrange/{\R{一}{いち}の くらい}/7/bkOrange} {
    \fill[\c!12, rounded corners=2mm] (\px,0) rectangle ++(28,26);
    \node[font=\sffamily\bfseries\fontsize{8.5}{9}\selectfont, text=\c!80!black] at (\px+14,22.8) {\t};
    \node[font=\sffamily\bfseries\fontsize{18}{18}\selectfont, text=\tc!80!black] at (\px+14,3.6) {\d};}
  \bkbill{(12.5,16.3)}{1000}\bkbill{(15.5,12.6)}{1000}
  \foreach \px in {30,60} {\node[font=\sffamily\bfseries\fontsize{9}{10}\selectfont, text=bkGray] at (\px+14,13.5) {なし};}
  \foreach \x in {97,104,111,118} {\bkcoin{(\x,16.5)}{1}}
  \foreach \x in {100.5,107.5,114.5} {\bkcoin{(\x,10)}{1}}
\end{tikzpicture}
\end{bkfigure}
\kotae 2007\R{円}{えん}（「にせん なな」と よむ）
\end{reidai}

\begin{jibun}
1000\R{円}{えん}さつ 1まいと、 10\R{円}{えん}\R{玉}{だま} 4まい。 なん\R{円}{えん}？
\bkkotaeran[6\zw]
\bkjans{1040\R{円}{えん}}
\end{jibun}

\begin{column}{0の はじまり}
なにも ない ことを、 「0」と いう かずに した のは、 いまから 1400ねんほど まえの インドの ひとたちだと いわれて います。
\end{column}

\clearpage
%% ---------------- p.42 まちがえやすいところ ----------------
\bksection{tsumazuki}{まちがえやすい ところ}

\begin{tsumazuki}{1から はかって しまう}
\begin{bkfigure}
\begin{tikzpicture}[x=1mm, y=1mm]
  \begin{scope}
  \bkruler{(0,0)}{6.5}{8}
  \bktape{(6.5,4)}{39}{bkRoute4}{}
  \node[anchor=west] at (-6,13.5) {\bkbatsu\enspace{\small 7？}};
  \end{scope}
  \begin{scope}[shift={(66,0)}]
  \bkruler{(0,0)}{6.5}{8}
  \bktape{(0,4)}{39}{bkRoute4}{}
  \node[anchor=west] at (-6,13.5) {\bkmaru\enspace{\small 6}};
  \end{scope}
\end{tikzpicture}
\end{bkfigure}
1に あわせると、 おわりの 7は ながさでは ありません。 0に あわせて よもう。
\end{tsumazuki}

\begin{tsumazuki}{ものさしの はしに あわせる}
ものさしの はしは 0では ない ことが あります。 \R{目}{め}もりの 0を さがして あわせよう。
\end{tsumazuki}

\begin{tsumazuki}{「ひゃくご」の かきかた}
\bkbatsu\enspace 15\hspace{2\zw}\bkbatsu\enspace 1005\hspace{3\zw}\bkmaru\enspace 105\par
\R{百}{ひゃく}の くらいに 1、 \R{十}{じゅう}の くらいに 0、 \R{一}{いち}の くらいに 5。 すうじを 3つ ならべて かきます。
\end{tsumazuki}



\clearpage
%% ---------------- p.43 つかわれているところ ----------------
\begin{butsuri}
はかる きかいは、 まず 0に あわせて から つかいます。
\begin{bkfigure}
\begin{tikzpicture}[x=1mm, y=1mm]
  % たいじゅうけい
  \fill[white, rounded corners=3mm] (0,0) rectangle (38,44);
  \fill[bkGray!35, rounded corners=3mm] (4,8) rectangle (34,38);
  \bkhakari{(19,25)}{10}{50}{10}{50}{0}{}
  \node[font=\sffamily\bfseries\small] at (19,4) {たいじゅうけい};
  % ストップウォッチ
  \begin{scope}[shift={(42,0)}]
  \fill[white, rounded corners=3mm] (0,0) rectangle (38,44);
  \fill[bkGray!60] (16.5,35) rectangle (21.5,39);
  \filldraw[fill=bkBlue!70, draw=bkBlue!60!black, line width=0.8pt] (19,23) circle (12.5);
  \bkdisplay{(7,17.5)}{24}{0}
  \node[font=\sffamily\bfseries\small] at (19,4) {ストップウォッチ};
  \end{scope}
  % エレベーター
  \begin{scope}[shift={(84,0)}]
  \fill[white, rounded corners=3mm] (0,0) rectangle (38,44);
  \node[font=\sffamily\bfseries\fontsize{9}{10}\selectfont] at (10,39) {にほん};
  \node[font=\sffamily\bfseries\fontsize{9}{10}\selectfont] at (28,39) {ヨーロッパ};
  \fill[bkGray!20, rounded corners=1.5mm] (4,9) rectangle (16,35) (22,9) rectangle (34,35);
  \foreach \y/\t in {14/1,22/2,30/3} {\bkqcbutton{(10,\y)}{\t}{white}}
  \foreach \y/\t in {14/0,22/1,30/2} {\bkqcbutton{(28,\y)}{\t}{white}}
  \draw[bkTheme, line width=1.2pt] (28,14) circle (4.6);
  \draw[bkStem, line width=1.2pt] (1,11) -- (37,11);
  \node[font=\sffamily\bfseries\small] at (19,4) {エレベーター};
  \end{scope}
\end{tikzpicture}
\end{bkfigure}
\begin{itemize}[label={\color{bkPurple}●}, leftmargin=1.4\zw, itemsep=0pt, topsep=0pt]
\item たいじゅうけいや ストップウォッチは、 はかる まえに 0に あわせます。
\item ヨーロッパには、 じめんの かいを 「0かい」と する くにが あります。
\item ロボットの うでも、 うごく まえに 「はじめの いち（0）」に もどってから うごきだします。
\end{itemize}
\end{butsuri}

\bkfukidashi{delta}{どこを 0に するか きめる。 それが 「やじるしの \R{路線}{ろ|せん}」の はじまりだよ！}

\begin{bkfigure}
\begin{tikzpicture}[x=1mm, y=1mm]
  \draw[bkRoute6, line width=3pt] (10,0) -- (128,0);
  \node[anchor=west, inner sep=0pt] at (0,12) {\bkrouteicon[1.4em]{6}\hspace{0.3em}{\sffamily\bfseries\small\color{bkRoute6!80!black}やじるしの \R{路線}{ろ|せん}}};
  \foreach \x/\c/\t/\s in {10/bkVol1/{1-3}/{0},39.5/bkVol1/{1-6}/{\R{数直線}{すう|ちょく|せん}},69/bkVol2/{2-2}/{マイナス},98.5/bkVol3/{3-2}/{ざひょう},128/bkVol4/{4-10}/{ベクトル}} {
    \fill[white, draw=\c, line width=1.6pt] (\x,0) circle (2.6);
    \node[anchor=base, font=\sffamily\bfseries\footnotesize, text=\c!85!black] at (\x,-8.4) {\t};
    \node[anchor=base, font=\sffamily\footnotesize, text=bkInk] at (\x,-15.8) {\s};}
  \node[fill=bkTheme, text=white, rounded corners=1mm, font=\sffamily\bfseries\fontsize{8}{9}\selectfont, inner sep=1.5pt] at (10,5.5) {いま ここ};
\end{tikzpicture}
\end{bkfigure}

\clearpage
%% ---------------- p.44〜45 もんだい ----------------
\bkendproblems

\bkproblemgroup{きほん}{1}

\bkproblem なんこ あるかな？ すうじで かこう。
\par\nobreak\vspace{2pt}
\begin{bkfigure}
\begin{tikzpicture}[x=1mm, y=1mm]
  \foreach \px/\l in {0/（1）,43/（2）,86/（3）} {\fill[white, draw=bkLine, line width=0.8pt, rounded corners=2.5mm] (\px,0) rectangle ++(40,30);
    \node[anchor=north west, font=\small, text=bkInk] at (\px+0.5,29.5) {\l};
    \node[inner sep=0pt] at (\px+20,-5) {\bkkakikomi\enspace{\small こ}};}
  \bkqcbasket{(20,16)}{5}{bkBlue}
  \bkqcbasket{(63,16)}{0}{bkBlue}
  \bkplate{(106,10)}{30}
  \foreach \x in {99,106,113} {\bkapple{(\x,13)}{3.2}{bkAppleRed}}
\end{tikzpicture}
\end{bkfigure}
\bkpans{（1）5\quad（2）0\quad（3）3}

\bkproblem あいて いる ますに はいる かずを かこう。\par\vspace{2pt}
\noindent\hspace*{1\zw}（1）4、\enspace 3、\enspace 2、\enspace\bkkakikomi、\enspace 0\hspace{3\zw}（2）3、\enspace 2、\enspace 1、\enspace\bkkakikomi
\bkpans{（1）1\quad（2）0}

\bkproblem ながさは、 \R{目}{め}もり いくつ\R{分}{ぶん}？
\par\nobreak\vspace{2pt}
\begin{bkfigure}
\begin{tikzpicture}[x=1mm, y=1mm]
  \node[anchor=west, font=\small] at (-6,14) {（1）};
  \bkruler{(4,0)}{6}{9}
  \bktape{(4,4)}{30}{bkRoute6}{}
  \node[font=\small] at (31,-15) {\bkkakikomi\enspace つ\R{分}{ぶん}};
  \begin{scope}[shift={(66,0)}]
  \node[anchor=west, font=\small] at (-6,14) {（2）};
  \bkruler{(4,0)}{6}{9}
  \bkpencil{(6.2,7.5)}{51.8}{bkRoute3}
  \node[font=\small] at (31,-15) {\bkkakikomi\enspace つ\R{分}{ぶん}};
  \end{scope}
\end{tikzpicture}
\end{bkfigure}
\bkpans{（1）5つ\R{分}{ぶん}\quad（2）9つ\R{分}{ぶん}}

\bkproblem かずを かこう。
\begin{enumerate}[itemsep=4pt, topsep=2pt]
\item 10の ぼうが 3\R{本}{ぼん}、 ばらは なし\enspace →\enspace\bkkakikomi
\item 100の いたが 2まい、 ばらが 4こ\enspace →\enspace\bkkakikomi
\end{enumerate}
\bkpans{（1）30\quad（2）204}

\bkproblemgroup{チャレンジ}{1}

\bkproblem おさらに リンゴが 3こ。 1こずつ たべて いくと、 おさらの リンゴは 3こ、 2こ、 1こ、\enspace\bkkakikomi\enspace こ。
\bkpans{0こ}

\bkproblem なん\R{円}{えん}かな？
\par\nobreak\vspace{2pt}
\begin{enumerate}[itemsep=4pt, topsep=2pt]
\item 100\R{円}{えん}\R{玉}{だま} 1まい、 10\R{円}{えん}\R{玉}{だま} 2まい\enspace →\enspace\bkkakikomi[15mm]\enspace\R{円}{えん}
\item 1000\R{円}{えん}さつ 1まい、 1\R{円}{えん}\R{玉}{だま} 5まい\enspace →\enspace\bkkakikomi[15mm]\enspace\R{円}{えん}
\end{enumerate}
\bkpans{（1）120\R{円}{えん}\quad（2）1005\R{円}{えん}}
\par\vspace{6mm}
\begin{furikaeri}
\begin{checklist}
\item なにも ない ことを 0で あらわせる。\hfill{\footnotesize\sffamily\color{bkTheme}→\pageref{bk:1-3:nai}ページ}
\item 0に あわせて ながさを よめる。\hfill{\footnotesize\sffamily\color{bkTheme}→\pageref{bk:1-3:awaseru}ページ}
\item なにも ない くらいに 0を かける。\hfill{\footnotesize\sffamily\color{bkTheme}→\pageref{bk:1-3:kurai}ページ}
\end{checklist}
チェックが つかなかったら、 →の ページに もどろう。
\end{furikaeri}

\bkfukidashi{sora}{おうちの ものさしや たいじゅうけいで、 0が どこに あるか さがして みよう！}
\clearpage
